\newpage % Nova página
\thispagestyle{empty} % Remove o número da página
\clearpage % Finaliza a página atual

\begin{center}
    \section*{Resumo} % Centraliza o título "Resumo" (sem numeração)
\end{center}

\par Os sistemas supervisórios (\emph{Supervisory Control and Data Acquisition}, SCADA) monitoram, supervisionam e controlam processos industriais em tempo real, reunindo em uma única interface a leitura das variáveis de campo, o registro histórico, a sinalização de alarmes e o envio de comandos aos equipamentos. Os sistemas convencionais, porém, exigem licenças proprietárias, hardware dedicado e instalação local especializada, o que os mantém fora do alcance de pequenas e médias empresas. Diante disso, este trabalho apresenta o desenvolvimento e a validação de um sistema supervisório web, multi-cliente e orientado a serviços, executado em contêineres Docker, acessado pelo navegador sem instalação no cliente. A arquitetura separa a interface web, a interface de programação de aplicações, os \emph{workers} de processamento em segundo plano, o serviço de aquisição Modbus TCP e um motor isolado para a execução de código do usuário em JavaScript. O banco de dados relacional PostgreSQL guarda o cadastro e o histórico, com um banco próprio para cada cliente, e o Redis mantém o valor atual de cada ponto e as filas de trabalho, enquanto os serviços se comunicam por um barramento de mensagens (Moleculer). O acesso é controlado por cargo e por grupo de permissões, e o comando de equipamentos segue quatro fases, de seleção, verificação, confirmação e execução, com trava por ponto, idempotência e registro de auditoria. São revisados os conceitos de sistemas supervisórios, Indústria 4.0, contêineres e orquestração que fundamentam a proposta, e os requisitos do sistema são validados contra a arquitetura. A validação combinou testes automatizados, campanhas de carga e uma bancada com controlador lógico programável TM221 e inversor de frequência ATV320. Os resultados medidos mostram a criação automática do banco de um cliente em até 3,4 s, a gravação do histórico de cerca de 9.400 linhas por segundo para um cliente e a aquisição de até 9.750 pontos a cada 30 s por réplica do serviço de aquisição, além de 8 ms para aplicar um ciclo de valores no sinótico. Os requisitos de comando, de alarmes e de replicação ainda dependem de ensaios na bancada e de medidas com mais de uma réplica, e a implantação em provedor de nuvem e em orquestrador de contêineres fica como trabalho futuro.

\leavevmode\newline
\par\textbf{Palavras-chave:} SCADA; supervisão de processos industriais; monitoramento remoto; contêineres; arquitetura orientada a serviços; Indústria 4.0; Modbus. % Negrito para "Palavras-chave", seguido das palavras-chave do trabalho

\thispagestyle{empty} % Remove a numeração da página
\clearpage % Nova página após o resumo
